% Section VI: Discussion
\section{Discussion}
\label{sec:discusion}

\subsection{Interpretation of the Results}
\label{sec:interpretacion}
The results separate two ways of understanding heterogeneity in attention. The first, combining two similarity functions within a head, does not improve on the better of its branches alone on any of the three datasets, and on CIFAR-10 it does not improve on dot-product attention once capacity is controlled. There are two reasons. When the combination happens before the \emph{softmax}, the mixture is a product of distributions and the branch with sharper scores decides; the relative scale of the scores acts as a temperature that picks the winner, and the other branch becomes uniform or irrelevant (Table~\ref{tab:colapso}). When the combination happens after the \emph{softmax}, or with standardized scores, both branches survive, but accuracy does not rise. Both branches answer the same question, which tokens resemble token $i$, from the same queries and keys; there is no complementary information to combine. The only fusion with a consistent gain, \var{ham\_concat}, combines \var{dot} with a global summary that is not a pairwise similarity function, and even so its accuracy advantage cannot be distinguished from that of a wider MLP.

The second way, adding to the head a term with an inductive bias that attention lacks, does work, and by a margin ten times larger. The $1\times1$ kernel control shows that what matters is spatial neighborhood: the same additional path from the values to the output, without mixing neighboring patches, changes nothing. The attention of a small ViT trained from scratch on 45,000 images is global in every layer (mean distance of 14 to 17~pixels on 32-pixel images) and does not learn by itself to mix neighbors efficiently. With the local term, the convolution takes over that mixing and attention reorganizes: it becomes more uniform and less diverse across heads in the intermediate layers, the layers in which \var{dot} was most selective. It is a division of labor between inductive biases, and it is the form of heterogeneity that helps in this regime.

The scope experiments show that this reading does not depend on the dataset, the augmentation or the model size within the tested range. The neighborhood separates from capacity on four datasets, holds with CutMix and with a strong augmentation, and its gain does not shrink between 119,000 and 3.6 million parameters. CutMix, which induces local heads in a 2.4-million-parameter ViT \cite{santiago2026inducing}, does not do so here. Nor does it matter much where the neighborhood is introduced, as long as it is in every block: in the attention output or in the MLP (LocalViT) it performs equally, whereas locality only at the input (convolutional stem, SPT) is much less effective. The comparison with a ResNet of the same size puts the result in perspective: at this scale, the model with the most convolutional bias is the best, and the local term brings the ViT closer to the CNN without reaching it.

The \emph{pooling} branch can be read the same way: a third inductive bias, a global summary of the image. On top of the local term it adds $+0.67$ points with 16 seeds, an advantage that does not survive correction for the planned contrasts; against an equal-capacity control it is $+0.49$ with 16 seeds and $+0.61$ with 48, significant in a contrast declared in advance. It is real, but an order of magnitude smaller than that of the neighborhood, and it raises the cost from 1.08 to 1.78 times the epoch time of \var{dot}. Pairwise attention is not indispensable either once there is a neighborhood: without it, with the local term and the \emph{pooling} branch, the model is $0.47$ points below \var{dot\_local\_k5} and $5.17$ above \var{dot}. The triad is also better calibrated than its control; part of that improvement may come from the per-head projection of \var{concat\_linear}, which without neighborhood explains all of the calibration improvement of \var{ham\_concat}.

Two further observations qualify the numbers. First, the variants with a local term have two training modes, and the larger neighborhood makes the high mode more likely (9 of 16 seeds with $5\times5$, against 2 of 16 with $3\times3$); with the \emph{pooling} branch it reaches 12 of 16 (31 of 48), against 8 of 16 (21 of 48) for its capacity control ($p=0.065$, Fisher). Part of the gain of $5\times5$ over $3\times3$ is that higher probability, and part is a shift of the whole distribution. Second, the local term improves calibration and robustness on CIFAR-10, but on PlantVillage it worsens robustness to patch occlusion; a bias towards the neighborhood makes the model more dependent on the neighbors being present.

PlantVillage refines the thesis. There, pairwise additive attention does improve on \var{dot}, and its effect adds almost entirely to that of the local term ($+1.00$ together, against $+0.52$ and $+0.63$ separately). On that dataset the additive branch behaves as a bias distinct from \var{dot} and not as a second version of the same similarity. Whether two attention mechanisms are complementary therefore depends on the data, and it is worth checking with the branch alone before combining it, because the fusion adds nothing over the dominant branch.

\subsection{What the Heterogeneous Architecture Contributes}
\label{sec:aporte}
Table~\ref{tab:evidencia} separates what each component of the triad contributes according to the strength of the evidence. Three contributions are supported, one of them small, and a fourth depends on the data.

\begin{itemize}
  \item \textbf{Confirmed: the neighborhood.} It is the only component with a large accuracy gain separated from capacity ($+5.31$ over \var{dot\_wide}, 16 of 16 seeds).
  \item \textbf{Confirmed: a stable way of combining branches.} Mixing scores before the \emph{softmax} collapses onto one branch; concatenating the outputs keeps both (the \emph{pooling} branch provides 37 to 49~\% of the fusion output) and is the only fusion with a positive sign on every dataset where it was evaluated: $+0.48$ on CIFAR-10, $+0.40$ on Food-101 and $+0.31$ on PlantVillage, only the first significant. The per-head projection of that fusion improves calibration (ECE $-0.61$), but the same effect appears with two identical \var{dot} branches: it is a contribution of the fusion, not of heterogeneity.
  \item \textbf{Supported by the primary contrast, but small: the global summary branch on top of the neighborhood.} The contrast declared before running seeds 16 to 47 supports it: the triad beats its equal-capacity control by $+0.61$ [$+0.17$, $+1.06$] with 48 seeds ($p=0.008$) and by $+0.67$ on the 32 new seeds ($p=0.025$), and is better calibrated (ECE $-0.36$, Holm $p=0.002$). The gain comes from entering the high mode more often (31 of 48 seeds against 21 of 48) and from raising the low mode. Without neighborhood, the branch improves on its controls in the same direction, but without surviving correction ($+0.37$ over the control with the same fusion and $+0.15$ over the wide MLP), and on PlantVillage it repeats the sign in all 5 seeds ($+0.24$ without neighborhood and $+0.23$ with it). In the scope experiments it beats its control on average in eight of nine conditions and ties in the ninth (Section~\ref{sec:res-triada-condiciones}), an exploratory check. The effect is an order of magnitude smaller than that of the neighborhood, and the triad costs 1.78 times the epoch time of \var{dot} (1.62 times that of its control), so it does not meet the cost criterion set in advance. Since the extension to 48 seeds was decided after seeing the first 16, the cleanest evidence is the replication on the new seeds.
  \item \textbf{Data-dependent: a second pairwise similarity.} On PlantVillage, pairwise additive attention adds to the neighborhood ($+0.38$ over \var{dot\_local\_k5}); on CIFAR-10, a second similarity on top of the neighborhood does not help ($+0.09$ with 5 seeds). In both cases the contribution comes from the branch and not from the fusion.
\end{itemize}

Overall, heterogeneity does help, but not where the original hypothesis looked for it. Combining two similarities over the same queries and keys does not help. What helps is a neighborhood bias that attention does not learn from little data and, to a much smaller extent, a global summary combined without mixing scores, whose contribution only separates from capacity on top of the neighborhood and with 48 seeds. \var{ham\_concat} should not be presented as an attention superior to \var{dot}, but as a way to integrate a second bias without collapse and with better calibration, although it costs 1.68 times the epoch time of \var{dot}.

% Translated by scripts/tables_en.py from the Spanish table generated by scripts/ieee_tables.py a partir de los runs; no editar a mano.

\begin{table*}[!t]
\centering
\footnotesize
\caption{What each component of the triad \texttt{ham\_concat\_local\_k5} contributes. Paired differences by seed on CIFAR-10 (top) and PlantVillage (bottom, batch 128, seeds 0--4); for ECE (\%), $\Delta>0$ counts the seeds in which it decreases. $n$: paired seeds. Holm over the rows of the table except the primary contrast (the triad against its capacity control in accuracy), which was declared as a single test before running seeds 16--47; the replication uses only those seeds. Suggestive: uncorrected $p<0.05$.}
\label{tab:evidencia}
\begin{adjustbox}{max width=\linewidth}
\begin{tabular}{>{\raggedright\arraybackslash}p{4.4cm}lccccc>{\raggedright\arraybackslash}p{1.9cm}}
\toprule
Component & Contrast & Metric & $\Delta$ [95~\% CI] & $\Delta>0$ / $n$ & $p$ & $p$ Holm & Evidence \\
\midrule
$5\times5$ neighborhood, at equal capacity & \texttt{dot\_local\_k5} $-$ \texttt{dot\_wide} & Accuracy & $+5.31$ [$+4.78$, $+5.83$] & 16/16 & $<0.001$ & $<0.001$ & confirmed \\
Per-head fusion projection (calibration) & \texttt{ham\_dotdot} $-$ \texttt{dot} & ECE & $-0.61$ [$-0.85$, $-0.38$] & 14/16 & $<0.001$ & $<0.001$ & confirmed \\
\emph{Pooling} branch without neighborhood, against the projection & \texttt{ham\_concat} $-$ \texttt{ham\_dotdot} & Accuracy & $+0.37$ [$+0.01$, $+0.73$] & 11/16 & 0.045 & 0.225 & suggestive \\
\emph{Pooling} branch without neighborhood, at equal capacity & \texttt{ham\_concat} $-$ \texttt{dot\_wide} & Accuracy & $+0.15$ [$-0.31$, $+0.61$] & 8/16 & 0.491 & 0.981 & not distinguishable \\
\emph{Pooling} branch on top of the neighborhood & \texttt{ham\_concat\_local\_k5} $-$ \texttt{dot\_local\_k5} & Accuracy & $+0.67$ [$+0.07$, $+1.27$] & 13/16 & 0.032 & 0.190 & suggestive \\
Triad, at equal capacity (primary contrast) & \texttt{ham\_concat\_local\_k5} $-$ \texttt{dot\_wide\_local\_k5} & Accuracy & $+0.61$ [$+0.17$, $+1.06$] & 31/48 & 0.008 & --- & confirmed \\
Triad, at equal capacity (replication: seeds 16--47 only) & \texttt{ham\_concat\_local\_k5} $-$ \texttt{dot\_wide\_local\_k5} & Accuracy & $+0.67$ [$+0.09$, $+1.26$] & 21/32 & 0.025 & --- & confirmed \\
Triad, at equal capacity (calibration) & \texttt{ham\_concat\_local\_k5} $-$ \texttt{dot\_wide\_local\_k5} & ECE & $-0.36$ [$-0.55$, $-0.18$] & 37/48 & $<0.001$ & 0.002 & confirmed \\
Second pairwise similarity on top of the neighborhood & \texttt{ham\_zscore\_local\_k5} $-$ \texttt{dot\_local\_k5} & Accuracy & $+0.09$ [$-1.06$, $+1.24$] & 2/5 & 0.836 & 0.981 & not distinguishable \\
\midrule
Pairwise additive similarity on top of the neighborhood & \texttt{additive\_free\_local\_k5} $-$ \texttt{dot\_local\_k5} & Accuracy & $+0.38$ [$+0.13$, $+0.62$] & 5/5 & 0.013 & 0.093 & suggestive \\
Triad, at equal capacity & \texttt{ham\_concat\_local\_k5} $-$ \texttt{dot\_wide\_local\_k5} & Accuracy & $+0.23$ [$-0.05$, $+0.51$] & 5/5 & 0.089 & 0.266 & not distinguishable \\
\emph{Pooling} branch without neighborhood, at equal capacity & \texttt{ham\_concat} $-$ \texttt{dot\_wide} & Accuracy & $+0.24$ [$+0.01$, $+0.48$] & 5/5 & 0.046 & 0.225 & suggestive \\
\bottomrule
\end{tabular}
\end{adjustbox}
\end{table*}


\subsection{Comparison with Previous Work}
\label{sec:previos}
That convolutional biases improve ViTs trained on little data is consistent with a broad line of work: convolutions in the token projection (CvT \citep{wu2021cvt}), hybrid blocks (CMT \citep{guo2022cmt}), attention initialized as a convolution (ConViT \citep{dascoli2021convit}), local window attention \citep{liu2021swin}, conditional position encodings \citep{chu2023conditional} and techniques specific to small datasets \citep{lee2021vision}. The local term evaluated here is the LePE encoding of CSWin \citep{dong2022cswin}. The contribution of this work is not the term itself, but the controlled measurement: with the same architecture, the same seeds and capacity controls, the neighborhood contributes on the order of 4.5 to 5.6 points on CIFAR-10, while combining attention mechanisms contributes at most half a point. Measured at the same budget and with the same protocol (Table~\ref{tab:baselines}), the published mechanisms are ordered by how much local bias they impose in each block: the convolutional stem \citep{xiao2021early} and SPT+LSA \citep{lee2021vision}, which act on the input or the temperature, fall far below; ConViT, whose locality is only an initialization, lies in between, and LocalViT \citep{li2021localvit} ties with the local term at a higher cost. That a 121,390-parameter ResNet beats them all agrees with the observation that ViTs need more data or more scale than CNNs to compete \citep{dosovitskiy2021image}. The observation that the attention of a ViT trained on little data is global even in the first layers agrees with the analysis of Raghu \emph{et al.}~\cite{raghu2021vision}, according to which local attention in the lower layers only appears with enough data. Santiago Toledo and Martínez \cite{santiago2026inducing} found that CutMix, and not other augmentations, suffices for a 2.4-million-parameter ViT to develop local heads. In our model, with the same CutMix protocol, the most local head still attends at a mean distance of 11 or 12~pixels, close to the 16.5 of uniform attention, and the accuracy of \var{dot} does not change. Both results are compatible if the ability of attention to learn locality depends on model scale; that study uses a single seed and this one does not test intermediate sizes with CutMix.

The gate after attention, which improves stability and performance in large language models \citep{qiu2025gated}, does not help in this model: the learned gate stays close to 1 in layers 2 to 4. Its benefit plausibly depends on scale and on phenomena, such as attention sinks, that do not appear in a four-layer model.

Regarding our previous work \cite{martinez2025reconfiguration}, the 1.3-point improvement of HAM B is not reproduced with the current protocol. HAM B is a special case of \var{ham\_concat}: the same \emph{pooling} branch with the fusion projection fixed to the sum of the two outputs. With a separate validation set, a cosine schedule and 16 seeds, \var{ham\_concat} gains $+0.48$ over \var{dot} and $+0.15$ over its capacity control (Table~\ref{tab:contrastes}). The difference from that result is explained by how it was measured: the maximum of 400 evaluations on the test set overestimates the final accuracy, and with a single run per variant a 1.3-point difference is within reach of the between-seed variation: across 16 seeds, the final accuracy of \var{dot} ranges from 78.42 to 80.73~\%. Without data augmentation, with 5 seeds, \var{ham\_concat} obtains $+0.19$ (Table~\ref{tab:otros}). HAM A, the Hadamard product of the two outputs, was not evaluated here.

Finally, the attention maps reinforce, in vision, the warning of Jain and Wallace~\cite{jain2019attention} about using attention as an explanation: two trainings of the same model produce maps that differ as much as those of two variants (Supplementary Fig.~\ref{fig:mapas}), so an isolated map is no evidence that a mechanism attends to more relevant regions.

\subsection{Implications for Future Applications}
\label{sec:implicaciones}
\begin{itemize}
  \item \textbf{For small ViTs with little data}, a $5\times5$ local term in the output of each head is a cheap and effective modification: 4,992 parameters (4.2~\%; 3,744 effective), 3.4~\% more MACs and 1.08 times the epoch time for, over an equal-capacity control, $+5.3$ points on CIFAR-10, $+5.4$ on CIFAR-100, $+2.8$ on SVHN and $+0.6$ on PlantVillage, and 3.7 times less GPU time to reach 79~\% validation accuracy, close to the final value of \var{dot}. The gain holds up to 3.6 million parameters and with CutMix. Being initialized at zero, it can be added to an existing model without changing its initial behavior.
  \item \textbf{If only accuracy matters} with about 120,000 parameters and $32\times32$ images, a small CNN is still the best option in this regime; the local term is the cheapest way to bring a ViT closer to it when global attention is needed, for example to scale the model or to combine it with other modalities.
  \item \textbf{For the design of heterogeneous mechanisms}, it is better to combine different inductive biases (neighborhood, global context, pairwise similarity) than several similarity functions over the same projections, and to combine them in separate subspaces or as residual terms, not by mixing scores before the \emph{softmax}.
  \item \textbf{For evaluation}, the differences between attention variants in this regime are of the order of the between-seed variation. Capacity controls, at least 16 paired seeds (on the order of 50 for half-point effects, such as that of the global branch) and correction for multiple comparisons are needed; and attention maps must be compared against the between-seed variation of the same model.
\end{itemize}

\subsection{Threats to Validity}
\label{sec:limitaciones}
\begin{itemize}
  \item \textbf{Internal validity.}
  \begin{itemize}
    \item \emph{Shared hyperparameters.} Hyperparameters were not tuned per variant, nor for the comparison mechanisms or the ResNet. Some fusion or published mechanism might improve with another learning rate or regularization. Since there was no hyperparameter search, no variant gained an advantage from it.
    \item \emph{Reimplementations.} LocalViT, SPT+LSA, ConViT and the convolutional stem were reimplemented at the scale of this model and fitted to the same budget; they are not the original models. In particular, ConViT uses gated attention in all four layers.
    \item \emph{Cost.} The epoch time depends on the GPU and on the implementation: heads are computed in a loop and the pairwise additive branches have no optimized kernel, which penalizes the fusions and the \emph{pooling} branch more than the output terms. The MACs and memory of Table~\ref{tab:costos} do not depend on the hardware: they confirm that the local term is cheap (3.4~\% more MACs) and show that the extra cost of the triad comes from the implementation (it has the MACs of its control). MACs do not count element-wise operations. Criterion 4 was set in advance on epoch time and is kept that way.
  \end{itemize}
  \item \textbf{Construct validity.}
  \begin{itemize}
    \item \emph{Interpretability.} Rollout does not include the mixing of the local term and does not represent the concatenation fusion exactly. Attention distance measures the locality of attention, not that of the whole model.
    \item \emph{Robustness.} It is limited to synthetic perturbations.
    \item \emph{Bimodality.} The two training modes have no explanation, and the threshold that separates them was set after observing them. Comparisons of their frequency are exploratory: with thresholds between 85.0 and 86.0~\% their direction holds, but not their significance (Supplementary Table~\ref{tab:umbral}).
  \end{itemize}
  \item \textbf{External validity.}
  \begin{itemize}
    \item \emph{Scale and data.} The results cover four datasets of $32\times32$ and $64\times64$ images and models of 119,000 to 3.6 million parameters, always trained from scratch. Pretraining, large datasets and higher-resolution images, where attention can learn locality by itself, were not evaluated. Food-101 was only used with the fusions of the first family.
    \item \emph{Deployment.} Inference was only measured on a desktop CPU, in fp32 and with PyTorch; microcontrollers, NPUs, quantization and compilers were not evaluated, so the study does not show that these models are viable on edge hardware. Having few parameters is not enough: activation memory and latency depend on how each operation is implemented (Section~\ref{sec:res-costos}).
    \item \emph{Architecture ablations.} The effect of changing the number of heads, or of applying the local term only in some layers, was not studied.
  \end{itemize}
  \item \textbf{Statistical conclusion validity.}
  \begin{itemize}
    \item \emph{Variant selection.} The variants extended to 16 seeds were chosen by their results on seeds 0 to 4. Seeds 5 to 15 act as confirmation, and they reduced the advantage of \var{ham\_concat\_local\_k5} from $+1.32$ to $+0.67$ and that of the variant with branches split across heads from $+0.74$ to $+0.19$.
    \item \emph{Extension of the primary contrast.} The contribution of the \emph{pooling} branch over capacity was confirmed with 48 seeds, with an extension decided after seeing the first 16. The replication on the 32 new seeds ($p=0.025$) mitigates, but does not eliminate, that bias. Outside that contrast there are only consistent signs with 5 seeds per condition, which share seeds and, mostly, images. The effect size, 0.6 points, is smaller than the between-seed deviation.
    \item \emph{Few seeds in the scope experiments.} With 5 seeds (3 for model size), the scope experiments only detect effects of the size of the neighborhood effect. One-point differences, such as that of LocalViT or that of the triad in each condition, remain within the between-seed variation.
  \end{itemize}
\end{itemize}
